\setcounter{chapter}{12}\chapter{Probabilités conditionnelles et loi binomiale}\label{ch:13}

{\small\sffamily\color{teal}S36 (rappels), S37, S38, S39, S41\par}

\begin{retenir}{Objectifs}Construire un arbre ; choisir une loi binomiale ; calculer une probabilité cumulée et un seuil ; interpréter espérance et variance.\end{retenir}

\section{Notes de cours}
\subsection{Rappels : conditionnement et indépendance}
Si $P(A)>0$, $P_A(B)=\frac{P(A\cap B)}{P(A)}$. Une branche d'arbre porte une probabilité conditionnelle ; la probabilité d'un chemin est le produit de ses branches. Pour une partition $A_1,\ldots,A_r$, $P(B)=\sum_iP(A_i\cap B)$ ; si tous les $P(A_i)>0$, elle s'écrit $\sum_iP(A_i)P_{A_i}(B)$.
Deux événements sont indépendants lorsque $P(A\cap B)=P(A)P(B)$. Pour $P(A)>0$, cela équivaut à $P_A(B)=P(B)$. Ne pas confondre $P_A(B)$ et $P_B(A)$.
\paragraph{Exemple.} $P(A)=0{,}4$, $P_A(B)=0{,}7$, $P_{\overline A}(B)=0{,}2$. Alors $P(B)=0{,}4\times0{,}7+0{,}6\times0{,}2=0{,}4$ et $P_B(A)=\frac{0{,}28}{0{,}4}=0{,}7$.
\subsection{Bernoulli et binomiale}
Une épreuve de Bernoulli a deux issues, succès de probabilité $p$ et échec de probabilité $1-p$. La variable indicatrice $B$ vaut $1$ en cas de succès et $0$ sinon : $E(B)=p$, $V(B)=p(1-p)$.
Un schéma de Bernoulli est une répétition de $n$ épreuves identiques et indépendantes. Le nombre $X$ de succès suit la loi binomiale $\mathcal B(n,p)$.
\begin{retenir}{Loi et paramètres}
Pour $k\in\{0,\ldots,n\}$,
\[P(X=k)=\binom nk p^k(1-p)^{n-k}.\]\[E(X)=np,\qquad V(X)=np(1-p),\qquad\sigma(X)=\sqrt{np(1-p)}.\]
\end{retenir}
\paragraph{Démonstration de la loi.} Chaque chemin comportant $k$ succès a probabilité $p^k(1-p)^{n-k}$ par indépendance. Il existe $\binom nk$ choix des positions des succès. Ces chemins sont incompatibles ; additionner leurs probabilités.
\subsection{Événements cumulés}
\[P(X\le k)=\sum_{j=0}^kP(X=j),\qquad P(X\ge k)=1-P(X\le k-1).\]\[P(a\le X\le b)=P(X\le b)-P(X\le a-1).\]
Pour au moins un succès, $P(X\ge1)=1-(1-p)^n$. Le complément de « au moins $k$ » est « strictement moins de $k$ ».
\paragraph{Exemple.} Pour $X\sim\mathcal B(10,0{,}2)$, $E(X)=2$ et $V(X)=1{,}6$. $P(X=2)=45(0{,}2)^2(0{,}8)^8\simeq0{,}3020$ et $P(X\ge1)=1-0{,}8^{10}\simeq0{,}8926$. L'espérance n'est ni un résultat certain ni nécessairement le plus probable.
\subsection{Seuil et intervalle}
Pour chercher le plus petit $n$ tel que $1-(1-p)^n\ge q$, une boucle sur $n$ suffit. Après le logarithme, on peut résoudre $(1-p)^n\le1-q$ puis arrondir dans le bon sens.
Pour trouver le plus petit entier $k$ tel que $P(X>k)\le\alpha$, augmenter $k$ jusqu'à obtenir $P(X\le k)\ge1-\alpha$, puis vérifier le rang précédent. On peut également chercher un intervalle $[a,b]$ contenant au moins une probabilité prescrite, sans supposer qu'il soit unique.
\subsection{Épreuves indépendantes, au-delà de deux issues}
Pour des expériences indépendantes, la probabilité d'une liste de résultats est le produit des probabilités de ses composantes. L'univers est le produit cartésien des univers de chaque expérience. Par exemple, pour deux lancers indépendants d'un dé équilibré, l'univers contient les $36$ couples $(i,j)$, $1\le i,j\le6$, chacun de probabilité $\frac{1}{36}$. La probabilité que la somme soit $5$ vaut $\frac{4}{36}=\frac{1}{9}$, car les couples favorables sont $(1,4),(2,3),(3,2),(4,1)$.
Dans un schéma de Bernoulli, on code le succès par $1$ et l'échec par $0$ : l'univers est $\{0,1\}^n$. L'indépendance concerne toute la succession ; une simple égalité entre deux probabilités observées ne suffit pas à la prouver.


\section{Activité de recherche}\textbf{Contrôle de composants --- 40 min.}\par On choisit au hasard un composant dans une production issue de deux ateliers. On note $A$ l'événement « le composant provient du premier atelier », $B=\overline A$ et $D$ l'événement « le composant est défectueux ». On donne $P(A)=0{,}6$, $P_A(D)=0{,}02$ et $P_B(D)=0{,}05$. Les résultats approchés seront arrondis à $10^{-4}$.
\begin{enumerate}\item Représenter la situation par un arbre pondéré, puis calculer $P(D)$.
\item Le composant choisi est défectueux. Calculer la probabilité qu'il provienne du premier atelier.
\item On choisit maintenant vingt composants indépendamment, chacun suivant le modèle précédent. On note $X$ le nombre de composants défectueux parmi les vingt. Justifier la loi suivie par $X$ et préciser ses paramètres.
\item Calculer $P(X=0)$, $P(X\ge2)$ et $E(X)$.
\item Si les vingt composants étaient prélevés dans un même lot, quelle hypothèse du modèle faudrait-il réexaminer ? Expliquer pourquoi.\end{enumerate}

\section{Exercices progressifs}

\exercice{1}{Entraînement}On considère deux événements $A$ et $B$ tels que $P(A)=0{,}3$, $P_A(B)=0{,}8$ et $P_{\overline A}(B)=0{,}1$. Calculer $P(B)$, puis $P_B(A)$.

\exercice{2}{Entraînement}La variable aléatoire $X$ suit la loi binomiale de paramètres $8$ et $\frac14$. Donner les valeurs exactes de $P(X=3)$, $E(X)$ et $V(X)$.

\exercice{3}{Entraînement}Soit $X$ une variable aléatoire à valeurs entières. Pour tout entier $k$, on pose $F(k)=P(X\le k)$. Exprimer $P(X\ge4)$ et $P(2\le X\le5)$ à l'aide de $F$.

\exercice{4}{Entraînement}On tire successivement cinq cartes, sans remise, dans un jeu de 52 cartes contenant quatre as. On note $X$ le nombre d'as obtenus. La variable $X$ suit-elle une loi binomiale ? Justifier.

\exercice{5}{Entraînement}On répète indépendamment une épreuve dont la probabilité de succès est $\frac12$. Déterminer le plus petit nombre d'épreuves pour lequel la probabilité d'obtenir au moins un succès est supérieure ou égale à $0{,}99$.

\section{Travail dans le manuel ou une annale}\textbf{Baccalauréat, Métropole, 17 juin 2025 : exercice 1, question 5 (p. 3).} En classe après S41. Lire le contexte des questions 1 à 4 ; pour ce travail, admettre le résultat $p=0{,}0714$ de la question 4. Le sujet est accessible ici : \href{https://www.education.gouv.fr/sites/default/files/2025-06/baccalaur-at-g-n-ral-2025-math-matiques-440856.pdf}{annale officielle}. Utiliser les données et les consignes de cette annale.

\section{Devoir à la maison}\textbf{DM13 --- Réservations et surréservation --- 60 min.}\par Une salle dispose de dix-huit places. L'organisateur accepte vingt réservations, payées à l'avance. Chaque personne ayant réservé se présente avec une probabilité de $0{,}8$, indépendamment des autres. On note $X$ le nombre de personnes qui se présentent. Si plus de dix-huit personnes viennent, les personnes excédentaires sont refusées.
\begin{enumerate}\item Justifier la loi de $X$, puis calculer son espérance et sa variance.
\item Calculer la probabilité que l'organisateur doive refuser au moins une personne. Arrondir à $10^{-6}$.
\item Chaque réservation rapporte $10$ unités monétaires, y compris si la personne ne vient pas. Chaque personne refusée reçoit un versement total de $30$ unités monétaires, remboursement et indemnité compris. On note $R=\max(X-18;0)$ le nombre de personnes refusées. Exprimer le gain $G$ en fonction de $R$, puis calculer $E(G)$.
\item Calculer le gain si l'organisateur n'accepte que dix-huit réservations, avec les mêmes règles de paiement. Comparer les deux stratégies du point de vue du gain moyen.
\item Citer une hypothèse du modèle qui peut être discutable en pratique, puis un critère autre que le gain monétaire à prendre en compte dans cette décision.\end{enumerate}

\ifprof\clearpage\section{Conduite et corrigés}\begin{profbox}{Repères pour le professeur}Réactiver les conditionnelles en S36 avant tout schéma de Bernoulli. L'indépendance est une hypothèse du modèle, elle ne découle pas de la seule répétition. Pour un tirage sans remise, ne pas invoquer automatiquement une loi binomiale. Les formules d'espérance et de variance sont utilisées en S39 ; leur preuve par somme d'indicatrices est reprise au chapitre 14. Un seuil demandé en décembre doit pouvoir être obtenu par une boucle, le logarithme n'étant pas encore enseigné.\end{profbox}

\subsection{Corrigé de l'activité}\begin{enumerate}
\item L'arbre porte les branches $A$ et $B$, de probabilités $0{,}6$ et $0{,}4$. Depuis $A$, les branches $D,\overline D$ portent $0{,}02,0{,}98$ ; depuis $B$, elles portent $0{,}05,0{,}95$.
\begin{center}\begin{tikzpicture}[grow=right,level distance=29mm,level 1/.style={sibling distance=40mm},level 2/.style={sibling distance=20mm},edge from parent/.style={draw,navy},every node/.style={font=\small}]
\node {} child {node {$B$} child {node {$\overline D$} edge from parent node[below] {$0{,}95$}} child {node {$D$} edge from parent node[above] {$0{,}05$}} edge from parent node[below] {$0{,}4$}} child {node {$A$} child {node {$\overline D$} edge from parent node[below] {$0{,}98$}} child {node {$D$} edge from parent node[above] {$0{,}02$}} edge from parent node[above] {$0{,}6$}};
\end{tikzpicture}\end{center}
Par probabilités totales, $P(D)=P(A)P_A(D)+P(B)P_B(D)=0{,}012+0{,}020=0{,}032$.
\item $P(D)>0$ et $P(A\cap D)=0{,}6\times0{,}02=0{,}012$. Donc $P_D(A)=\frac{P(A\cap D)}{P(D)}=\frac{0{,}012}{0{,}032}=\frac{3}{8}=0{,}375$.
\item Chaque prélèvement est une épreuve à deux issues, « défectueux » ou « non défectueux », avec la même probabilité $0{,}032$ de défaut. Les vingt prélèvements étant indépendants, $X\sim\mathcal B(20,0{,}032)$.
\item $P(X=0)=0{,}968^{20}\simeq0{,}5218$. Par complément,
\[P(X\ge2)=1-P(X=0)-P(X=1)=1-0{,}968^{20}-20(0{,}032)0{,}968^{19}\simeq0{,}1332.\]
Enfin $E(X)=20\times0{,}032=0{,}64$. Cette espérance n'implique pas qu'un prélèvement présente un nombre fractionnaire de défauts.
\item Une cause commune liée au lot pourrait rendre les défauts dépendants. La loi binomiale et ses calculs reposent sur l'indépendance donnée dans le modèle ; il faudrait réexaminer cette hypothèse pour un autre protocole.
\end{enumerate}

\paragraph{E1.}Par probabilités totales, $P(B)=0{,}3\times0{,}8+0{,}7\times0{,}1=0{,}31$. Puisque $P(B)>0$, $P_B(A)=\frac{P(A\cap B)}{P(B)}=\frac{0{,}24}{0{,}31}=\frac{24}{31}$.

\paragraph{E2.}Par la formule binomiale, $P(X=3)=\binom{8}{3}(\frac{1}{4})^3(\frac{3}{4})^5=56(\frac{1}{4})^3(\frac{3}{4})^5$. Les paramètres donnent $E(X)=\frac{8}{4}=2$ et $V(X)=8(\frac{1}{4})(\frac{3}{4})=\frac{3}{2}$.

\paragraph{E3.}Le complément de $X\ge4$ est $X\le3$, donc $P(X\ge4)=1-F(3)$. Pour l'intervalle, on enlève à $X\le5$ les valeurs $X\le1$, donc $P(2\le X\le5)=F(5)-F(1)$.

\paragraph{E4.}Non, dans un jeu usuel contenant quatre as. Sans remise, obtenir un as au premier tirage change la probabilité d'en obtenir un au second : les essais ne sont pas indépendants. En outre, le nombre d'as ne peut jamais valoir $5$, alors qu'une binomiale de paramètres $5$ et $0<p<1$ donne une probabilité positive à cette valeur.

\paragraph{E5.}La probabilité d'au moins un succès vaut $1-(\frac{1}{2})^n$. La condition est équivalente à $2^n\ge100$. Comme $2^6=64<100\le128=2^7$, le premier entier est $7$. La probabilité vaut alors $\frac{127}{128}\simeq0{,}99219$.

\subsection{Corrigé du DM13}\begin{enumerate}
\item $X$ compte les succès de vingt épreuves de Bernoulli indépendantes de paramètre $0{,}8$. Donc $X\sim\mathcal B(20,0{,}8)$, $E(X)=20\times0{,}8=16$ et $V(X)=20\times0{,}8\times0{,}2=3{,}2$.
\item La salle manque de places si $X>18$, c'est-à-dire $X=19$ ou $X=20$. Ces événements étant incompatibles,
\[P(X>18)=20(0{,}8)^{19}(0{,}2)+(0{,}8)^{20}\simeq0{,}069175.\]
Le risque est d'environ $6{,}92\,\%$.
\item Le nombre de refus $R$ vaut $0$ pour $X\le18$, $1$ pour $X=19$ et $2$ pour $X=20$. Par définition de l'espérance,
\[E(R)=P(X=19)+2P(X=20)\simeq0{,}0807045.\]
Les recettes sont $20\times10=200$ et les versements $30R$, donc $G=200-30R$. Ainsi $E(G)=200-30E(R)\simeq197{,}5789$ unités monétaires. Les $30$ incluent déjà tous les versements liés à un refus ; il ne faut pas retirer encore $10$.
\item Avec dix-huit réservations, aucun refus n'est possible et les recettes sont $18\times10=180$, quel que soit le nombre de présents. L'espérance du gain avec vingt réservations dépasse ce montant d'environ $17{,}58$ unités monétaires.
\item Les venues de personnes d'un même groupe peuvent être dépendantes. En outre, la satisfaction du public et les conséquences d'un refus ne sont pas résumées par l'espérance du gain. Le résultat chiffre un modèle ; il ne suffit pas à lui seul à imposer une décision.
\end{enumerate}

\subsection{Corrigé du travail dans le manuel ou une annale}Les $100$ choix sont indépendants, avec une même probabilité $p=0{,}0714$. Donc $X\sim\mathcal B(100;0{,}0714)$.
\[P(X\le7)=\sum_{k=0}^7\binom{100}k(0{,}0714)^k(0{,}9286)^{100-k}\simeq0{,}577.\]
Enfin, $E(X)=7{,}14$ et $V(X)=100\times0{,}0714\times0{,}9286=6{,}630204\simeq6{,}63$.

\fi

